Learning-augmented (``with-predictions'') algorithms for online
bipartite matching have proliferated: MinPredictedDegree consumes a
per-node degree predictor, and a family of \emph{test-and-fallback}
schemes tests a type-histogram prediction on a sublinear prefix of the
arrivals before committing to follow it or to fall back on an
advice-free baseline. These algorithms have been analyzed in isolation
--- on different input models, against different error models, and
largely in theory. We give the first unified experimental study, placing
all of them on a single harness with a common structured
prediction-error model, a common optimum, and confidence intervals,
across synthetic graphs, six real-world graphs, and real request traces.
Our central finding is that, on average-case (known-i.i.d.) inputs, the
value of predictions is \textbf{robustness insurance, not a performance
lever}: the advice-free baseline is already near-optimal, so the
\emph{consistency} upside of good advice is small, whereas unguarded
prediction-following can crash far below the baseline --- and the
practical worth of the sophisticated algorithms is precisely that they
never do. We then prove what the wall \emph{costs}. For the
test-and-fallback class we establish a two-sided \textbf{budget--stakes
law}: deciding whether to follow the advice requires --- and, via an
explicit one-line rule, only requires --- a prefix of length
\(\tilde\Theta(\sigma^2/g^2)\): the instance's specialist mass
\(\sigma^2\) (exactly twice its baseline slack) divided by the square of
the stakes \(g\) (what good advice gains over the baseline), independent
of the instance length. The two halves meet, up to logarithms, whenever
the stakes are carried at comparable signal levels by a constant
fraction of that mass; one low-signal regime is left open and is stated
as such. The lower half binds \emph{every} decision rule; the upper half
refutes the natural conjecture that the hardness of tolerant
distribution testing blocks all sublinear rules --- the
decision-relevant statistic is the \emph{payoff} of following, which is
exponentially cheaper to test than the prediction's distance --- a
prescription we validate on the benchmark, where a constant-free
payoff-testing rule avoids the threshold pathologies of the deployed
tests. Because the stakes are capped by the baseline slack, on
strong-baseline instances the price exceeds the horizon: upsides below
\(\approx\sqrt{(1-\rho_{\mathrm{base}})/n}\) are uncapturable by any
rule at any prefix length, and the upsides we measure sit in that range.
Experiments and theory deliver one message: \textbf{on average-case
matching, the advice's upside is smaller than the price of finding out
whether to trust it.}
